\chapter{Group theory and topology}\label{ch:topology}

\begin{watch}{\lec{8} (no captions; times read off the slides)}
Band topology \ts{8}{600}{$\approx$10:00} $\cdot$
Wannier representability \ts{8}{1100}{$\approx$18:20} $\cdot$
Wannier representability and TQC \ts{8}{1500}{$\approx$25:00} $\cdot$
classification of band representations \ts{8}{2100}{$\approx$35:00} $\cdot$
EBRs from Wyckoff positions \ts{8}{2700}{$\approx$45:00} $\cdot$
irreps from EBRs \ts{8}{3300}{$\approx$55:00} $\cdot$
double-valued EBRs of $p6mm$ \ts{8}{3900}{$\approx$65:00} $\cdot$
spinful $p_z$ orbitals in graphene \ts{8}{4400}{$\approx$73:20} $\cdot$
EBR analysis \ts{8}{4800}{$\approx$80:00} $\cdot$
what TQC can and cannot diagnose \ts{8}{5150}{$\approx$85:50}
\end{watch}

Topological quantum chemistry (TQC) connects the symmetry labels of \cref{ch:bands}
with the global, topological properties of bands. Its central idea fits in one
sentence: \emph{a set of bands is topologically trivial exactly when it can be
described by localised, symmetric orbitals in real space; the irreps at
high-symmetry points often reveal when this is impossible.}

\section{Local versus global}

A cylinder and a M\"obius strip are made by gluing two opposite edges of a rectangle,
with or without a twist. Every small piece of either looks like a piece of rectangle:
locally they are identical. Globally they differ, and no continuous deformation turns
one into the other. Topology studies such global properties.

Bands carry the same kind of structure. The Bloch Hamiltonian is periodic in
reciprocal space, $\mat{H}(\vec{k}+\vec{G})=\mat{H}(\vec{k})$, so the Brillouin zone is
a torus. At each $\vec{k}$ an isolated set of $n$ bands defines an $n$-dimensional
subspace of states. As $\vec{k}$ runs over the torus these subspaces form a
\emph{vector bundle}. Locally one can always choose smooth basis vectors (a gauge).
Whether one can choose them smoothly and periodically over the \emph{whole} zone is a
global question, and the answer can be no. The Chern number of a two-dimensional band
is the simplest obstruction.

\section{Wannier functions}

From Bloch states $\psi_{n\vec{k}}$ of an isolated set of bands one builds
\emph{Wannier functions}
\begin{equation}
w_{n\vec{t}}(\vec{r})=\frac{V_{\mathrm{cell}}}{(2\cpi)^d}\int_{\mathrm{BZ}}\dd^d k\;
\ee^{-\ii\vec{k}\cdot\vec{t}}\sum_m U_{mn}(\vec{k})\,\psi_{m\vec{k}}(\vec{r}),
\end{equation}
labelled by a band index and a lattice vector, and related by translations,
$w_{n\vec{t}}(\vec{r})=w_{n\vec{0}}(\vec{r}-\vec{t})$. The unitary matrix
$\mat{U}(\vec{k})$ is the gauge freedom. Wannier functions are the real-space,
orbital-like description of the bands, and their localisation depends on the gauge:
a smooth, periodic gauge gives exponentially localised functions.

\begin{keyresult}[Wannier representability]
The Wannier functions of an isolated set of bands can be chosen
\emph{exponentially localised and transforming among themselves under all the
symmetries of the crystal} if and only if the bands are topologically trivial. Such a
set of bands is called \emph{Wannier representable}.
\end{keyresult}
In TQC this equivalence \emph{defines} ``trivial''. For two-dimensional insulators
without extra symmetry it contains the theorem that exponentially localised Wannier
functions exist precisely when all Chern numbers vanish. The global property in
reciprocal space (topology) is thus tied to a local property in real space
(localisation of symmetric orbitals).

\section{Band representations}

\subsection{Definition}
Take a Wyckoff position with sites $\vec{q}_1,\dots,\vec{q}_n$ in the cell,
site-symmetry group $G_{\vec{q}_1}$, and a set of orbitals at $\vec{q}_1$ transforming as
a representation $\rho$ of $G_{\vec{q}_1}$. Copy them to the other sites with the
space-group operations (coset representatives: the cosets of $G_{\vec{q}_1}$ in $G$
label the sites of the orbit, \cref{sec:cosets}) and to all cells with translations.
The infinite set of orbitals so obtained is mapped onto itself by the space group:
it carries a representation of $G$, the \emph{band representation}
\begin{equation}
\rho\uparrow G\qquad\text{(the representation \emph{induced} from $\rho$)}.
\end{equation}
In the Bilbao notation, $\rho@w$ ($w$ being the Wyckoff label); for instance
$\irr{A}_1@2b$ is our graphene $p_z$ example. Fourier transforming to Bloch sums, the
band representation becomes, at each $\vec{k}$, a finite representation of the little
group, of dimension $n\dim\rho$, whose character is exactly
\eqref{eq:bandchar}. Its small-irrep content at the high-symmetry points is
the ``irrep vector'' of the band representation.

A set of bands is Wannier representable if and only if it carries (is equivalent to)
a band representation. The Wannier functions are then the orbitals from which the
band representation was induced.

\subsection{Elementary band representations}
Band representations can be added. A band representation that cannot be written as a
sum of other band representations is \emph{elementary} (EBR). EBRs are the building
blocks of all atomic-limit band structures.
\begin{itemize}
\item Every EBR is induced from an irrep of the site-symmetry group of a \emph{maximal}
      Wyckoff position (one whose site-symmetry group is not contained in that of
      another position). Orbitals at non-maximal positions can be moved continuously
      to a maximal one without breaking symmetry, so they give composite band
      representations. (A few irreps at maximal positions also give composite band
      representations; these exceptions are listed in the tables.)
\item For each space group the number of EBRs is finite. They are tabulated, with
      their irrep vectors, in the \texttt{BANDREP} application of the Bilbao
      Crystallographic Server, for all 230 space groups (with and without spin--orbit
      coupling) and for the magnetic groups.
\end{itemize}

For $p6mm$ the maximal positions are $1a$ ($\Cv{6}$), $2b$ ($\Cv{3}$) and $3c$
($\Cv{2}$). Evaluating \eqref{eq:bandchar} for every irrep of each site group gives
the spinless EBRs:
\begin{equation}\label{eq:ebrp6mm}
\begin{array}{llccc}
\toprule
\rho@w & \text{bands} & \Gamma\ (\Cv{6}) & K\ (\Cv{3}) & M\ (\Cv{2})\\
\midrule
\irr{A}_1@1a & 1 & \irr{A}_1 & \irr{A}_1 & \irr{A}_1\\
\irr{A}_2@1a & 1 & \irr{A}_2 & \irr{A}_2 & \irr{A}_2\\
\irr{B}_1@1a & 1 & \irr{B}_1 & \irr{A}_2 & \irr{B}_1\\
\irr{B}_2@1a & 1 & \irr{B}_2 & \irr{A}_1 & \irr{B}_2\\
\irr{E}_1@1a & 2 & \irr{E}_1 & \irr{E} & \irr{B}_1\oplus\irr{B}_2\\
\irr{E}_2@1a & 2 & \irr{E}_2 & \irr{E} & \irr{A}_1\oplus\irr{A}_2\\
\midrule
\irr{A}_1@2b & 2 & \irr{A}_1\oplus\irr{B}_1 & \irr{E} & \irr{A}_1\oplus\irr{B}_1\\
\irr{A}_2@2b & 2 & \irr{A}_2\oplus\irr{B}_2 & \irr{E} & \irr{A}_2\oplus\irr{B}_2\\
\irr{E}@2b & 4 & \irr{E}_1\oplus\irr{E}_2 & \irr{A}_1\oplus\irr{A}_2\oplus\irr{E} &
   \irr{A}_1\oplus\irr{A}_2\oplus\irr{B}_1\oplus\irr{B}_2\\
\midrule
\irr{A}_1@3c & 3 & \irr{A}_1\oplus\irr{E}_2 & \irr{A}_1\oplus\irr{E} & \irr{A}_1\oplus\irr{B}_1\oplus\irr{B}_2\\
\irr{A}_2@3c & 3 & \irr{A}_2\oplus\irr{E}_2 & \irr{A}_2\oplus\irr{E} & \irr{A}_2\oplus\irr{B}_1\oplus\irr{B}_2\\
\irr{B}_1@3c & 3 & \irr{B}_1\oplus\irr{E}_1 & \irr{A}_2\oplus\irr{E} & \irr{A}_1\oplus\irr{A}_2\oplus\irr{B}_1\\
\irr{B}_2@3c & 3 & \irr{B}_2\oplus\irr{E}_1 & \irr{A}_1\oplus\irr{E} & \irr{A}_1\oplus\irr{A}_2\oplus\irr{B}_2\\
\bottomrule
\end{array}
\end{equation}
Conventions: the $\Cv{6}$ labels are those of \eqref{eq:C6vtable}. At $1a$ the site
irrep $\irr{B}_1$ is even under the bond-containing mirrors $\sv$. At $3c$ the site
irrep $\irr{B}_1$ is even under the mirror containing the bond, i.e.\ it is a
$p$ orbital pointing along the bond. At $M$, $\irr{B}_1$ is odd under $C_2$ and even
under $\sv$. The row $\irr{A}_1@2b$ is exactly the $\uppi$-band result of \cref{ch:bands}.
In the Bilbao labels the same information appears as vectors of multiplicities, one
entry per small irrep; always check the conventions of the table you use.

\section{The TQC criterion}\label{sec:tqc}

Every band representation has an irrep vector that is a sum of EBR vectors with
non-negative integer coefficients. Consequently:
\begin{keyresult}[Topological quantum chemistry]
Take an isolated set of bands (separated by gaps from the rest) and list its small
irreps at the maximal $\vec{k}$ points.
\begin{enumerate}
\item If this irrep vector \emph{cannot} be written as a sum of EBR vectors with
      non-negative integer coefficients, the set is \emph{necessarily} topological:
      no false positives.
\item If it can be written only with some negative coefficients, the set is
      \emph{fragile} topological: it becomes trivial when certain trivial bands are
      added.
\item If it can be written with non-negative coefficients, the set \emph{may} still
      be topological: false negatives exist, because not all topology is
      \emph{symmetry indicated}. Two sets of bands can have identical irreps
      everywhere and different topology (a Chern insulator and a trivial insulator,
      for example).
\end{enumerate}
\end{keyresult}
A second, purely structural statement is often decisive.
\begin{theorem}[Disconnected EBRs]
If the bands of an EBR split into two or more groups separated by gaps, at least one
of the groups is topological.
\end{theorem}
Indeed, if every group were Wannier representable, each would carry a band
representation, and the EBR would be their sum, contradicting its elementarity.

The compatibility relations of \cref{sec:compat} are what make band structures
connected or disconnected. Along each line the subduced irreps must match. If there is
no way of connecting the irreps at the high-symmetry points into two separated groups
that satisfy all compatibility relations, the EBR is forced to be connected. A
disconnection, when allowed, is a signal of topology.

\section{Graphene without spin--orbit coupling}

For the spinless $\uppi$ bands, $\irr{A}_1@2b$ has the two-dimensional irrep $\irr{E}$ at $K$.
The two bands are forced to touch there, so the EBR cannot split: graphene is a
semimetal with symmetry-enforced Dirac points. No gap, no topological insulator.

\section{Graphene with spin--orbit coupling}

With spin, each $p_z$ orbital becomes a Kramers doublet. At the $\Cv{3}$ site,
$\irr{a}_1\otimes D_{1/2}=\bar{\irr{E}}_1$ (\cref{ch:spin}), so the four spinful $\pi$
bands carry the double-valued EBR $\bar{\irr{E}}_1@2b$. Using the double-group
characters of $\Cv{6}$, $\Cv{3}$ and $\Cv{2}$ (\cref{ex:8:spinful}) one finds
\begin{equation}\label{eq:spinfulgraphene}
\bar{\irr{E}}_1@2b:\qquad
\Gamma:\ \bar{\irr{E}}_1\oplus\bar{\irr{E}}_2,\qquad
K:\ \bar{\irr{E}}_1\oplus{}^{1}\bar{\irr{E}}\oplus{}^{2}\bar{\irr{E}},\qquad
M:\ 2\,\bar{\irr{E}} .
\end{equation}
In the Bilbao labels shown in the lecture these are
$\Gamma_7\oplus\Gamma_8$, $K_4\oplus K_5\oplus K_6$ and $2M_5$, with $K_4$ the
two-dimensional irrep at $K$.

With spin--orbit coupling the degeneracy at $K$ need not be fourfold: the four bands can
separate into two isolated pairs. In the lecture's example (a Kane--Mele-type model),
\begin{itemize}
\item one pair carries $\{\Gamma_7\,|\,K_4\,|\,M_5\}$. This is exactly the irrep vector of
      the EBR $\bar{\irr{E}}_1@1a$ (spinful orbitals at the hexagon centres). By irreps
      alone it \emph{could} be trivial: an ``obstructed atomic limit'' with Wannier
      centres in the empty hexagons rather than on the atoms;
\item the other pair carries $\{\Gamma_8\,|\,K_5\oplus K_6\,|\,M_5\}$. The two-band EBRs
      of this group are $\bar{\irr{E}}_1@1a=\{\Gamma_7|K_4|M_5\}$,
      $\bar{\irr{E}}_2@1a=\{\Gamma_8|K_4|M_5\}$ and $\bar{\irr{E}}_3@1a=\{\Gamma_9|K_5\oplus K_6|M_5\}$.
      None of them matches, so this pair is necessarily topological.
\end{itemize}
This is consistent with the theorem on disconnected EBRs. In the Kane--Mele model the
topological pair is the $\mathbb{Z}_2$ quantum spin Hall insulator. There, in fact,
\emph{both} pairs are $\mathbb{Z}_2$ non-trivial. The pair whose irreps mimic an atomic
limit is a false negative of the irrep criterion, a concrete instance of point 3 above.

\begin{beyond}[Beyond the lectures: where to go next]
\begin{itemize}
\item The \texttt{BANDREP}, \texttt{Check Topological Mat.}\ and Topological Materials
      Database tools of the Bilbao server apply this analysis to first-principles band
      structures; programs such as \texttt{IrRep} extract the irreps from DFT output.
\item Symmetry indicators (Po, Vishwanath, Watanabe) are the group of irrep vectors
      modulo band representations; they classify what irreps can diagnose.
\item Magnetic topological quantum chemistry extends everything to the 1651 magnetic
      space groups, where time reversal combines with spatial operations
      (\cref{sec:magnetic}).
\item Review: J.\ Cano and B.\ Bradlyn, Annu.\ Rev.\ Condens.\ Matter Phys.\ 12,
      225 (2021).
\end{itemize}
\end{beyond}

\exercisesection

\begin{exercise}\label{ex:8:A1at1a}
Derive the rows $\irr{A}_1@1a$ and $\irr{E}_1@1a$ of \eqref{eq:ebrp6mm}. For orbitals
at the origin all phases in \eqref{eq:bandchar} are trivial. Why does $\irr{E}_1@1a$
split at $M$?
\end{exercise}

\begin{solution}
An orbital at the origin is mapped onto itself by every point operation, with
$\vec{t}_\alpha=0$, so $\chi_{\vec{k}}(g)=\chi_\rho(g)$: the band representation at $\vec{k}$ is
simply $\rho$ subduced to $\bar{G}_{\vec{k}}$. For $\irr{A}_1$ this is the identity irrep
everywhere. For $\irr{E}_1$: at $\Gamma$, $\irr{E}_1$; at $K$ the characters
$(2,-1,0)$ on $(E,C_3,\sd)$ give $\irr{E}$; at $M$ the characters $(2,-2,0,0)$ on
$(E,C_2,\sv,\sd)$ give $\irr{B}_1\oplus\irr{B}_2$. It splits because $\Cv{2}$ has only
one-dimensional irreps.
\end{solution}

\begin{exercise}\label{ex:8:notsum}
Show that $\irr{A}_1@2b$ is not the sum of EBRs induced from $1a$. Then explain why the
two spinless $\uppi$ bands can never be separated into two isolated bands without
breaking a symmetry of $p6mm$.
\end{exercise}

\begin{solution}
One-band EBRs from $1a$ have one-dimensional irreps at $K$, so any sum of two of them
has two one-dimensional irreps at $K$, not $\irr{E}$. The two-band EBRs from $1a$ have
$\irr{E}_1$ or $\irr{E}_2$ at $\Gamma$, not $\irr{A}_1\oplus\irr{B}_1$. Hence $\irr{A}_1@2b$ is not a
sum of $1a$ EBRs (the Wannier functions of the $\uppi$ bands sit on the atoms). The two
$\uppi$ bands can never be separated: at $K$ they form the two-dimensional irrep $\irr{E}$,
so they touch there for every symmetric Hamiltonian.
\end{solution}

\begin{exercise}\label{ex:8:obstructed}
A hypothetical isolated band of a $p6mm$ crystal has irreps $\irr{B}_1$ at $\Gamma$,
$\irr{A}_2$ at $K$ and $\irr{B}_1$ at $M$. Is it Wannier representable according to the
irrep criterion? Where would its Wannier centres be?
\end{exercise}

\begin{solution}
$\{\irr{B}_1|\irr{A}_2|\irr{B}_1\}$ is exactly the irrep vector of $\irr{B}_1@1a$
\eqref{eq:ebrp6mm}. By the irrep criterion the band is compatible with a Wannier
function centred at the hexagon centres ($1a$), with the symmetry of an
$f$-like orbital (odd under the bond-bisecting mirrors). If the atoms sit at $2b$, such a
band is an \emph{obstructed atomic limit}: trivial, but with Wannier centres away from
the atoms. Irreps cannot exclude hidden topology (false negatives), but nothing in them
requires it.
\end{solution}

\begin{exercise}[spinful $p_z$ orbitals]\label{ex:8:spinful}
The double group of $\Cv{6}$ has three two-dimensional double-valued irreps with
characters (classes $E$, $\Ebar$, $C_2$, $C_3$, $\bar{C}_3$, $C_6$, $\bar{C}_6$, $\sv$, $\sd$)
\[
\begin{array}{c|ccccccccc}
 & E & \Ebar & C_2 & C_3 & \bar{C}_3 & C_6 & \bar{C}_6 & \sv & \sd\\\hline
\bar{\irr{E}}_1 & 2&-2&0&1&-1&\sqrt3&-\sqrt3&0&0\\
\bar{\irr{E}}_2 & 2&-2&0&1&-1&-\sqrt3&\sqrt3&0&0\\
\bar{\irr{E}}_3 & 2&-2&0&-2&2&0&0&0&0
\end{array}
\]
(class sizes $1,1,2,2,2,2,2,6,6$; $C_2$ and $\bar{C}_2$, and each mirror and its barred
partner, fall in the same class).
Using \eqref{eq:bandchar} with the spin character $\chi_{1/2}(\theta)=2\cos(\theta/2)$
for rotations and $0$ for mirrors, derive \eqref{eq:spinfulgraphene}.
\end{exercise}

\begin{solution}
Each site carries $\irr{a}_1\otimes D_{1/2}$; its site character is the spin character.
At $\Gamma$: $E\to4$, $\Ebar\to-4$; $C_2$, $C_6$, $\sd$ exchange sublattices $\to0$; $C_3$
keeps them with trivial phases, $\to2\chi_{1/2}(120^\circ)=2$, and $\bar{C}_3\to-2$;
$\sv\to2\cdot0=0$. The reduction in the double group of $\Cv{6}$ (order 24, class sizes
$1,1,2,2,2,2,2,6,6$) gives
$m_{\bar{\irr{E}}_1}=m_{\bar{\irr{E}}_2}=\tfrac1{24}(8+8+4+4)=1$ and
$m_{\bar{\irr{E}}_3}=\tfrac1{24}(8+8-8-8)=0$.
At $K$: $C_3$ gives $(\omega+\omega^{*})\cdot1=-1$, $\bar{C}_3$ gives $+1$, the mirrors
$0$, so $\chi_K=(4,-4,-1,1,0,0)=\chi_{D_{3/2}}$, i.e.\
$\bar{\irr{E}}_1\oplus{}^{1}\bar{\irr{E}}\oplus{}^{2}\bar{\irr{E}}$ by \eqref{eq:C3vdouble}.
At $M$: $(4,-4,0,0,0)$ on the five classes of the double group of $\Cv{2}$ (order 8),
whose only double-valued irrep $\bar{\irr{E}}$ has $(2,-2,0,0,0)$: $2\,\bar{\irr{E}}$.
\end{solution}
